% Card: bridging (rook endings)
\paragraph{Bridging}\index{Rook endings!Bridging}

\diagram[bridging]{8/6P1/5K2/6R1/1k6/8/6r1/8 b - - 27 83}

Here the white rook went to g5 instead. It stands on the g-file between the black rook and the
pawn, like a bridge: the black rook can no longer use that file to stop the pawn. A check from the
side changes nothing: after 83...\R f2+ 84.\K g6 the pawn queens. This is
\vocab[rook endings]{bridging}{bridging}{The attacker puts his rook between the defending rook and
his own king and pawn, so that the checks stop and the pawn can queen}.
